\chapter{Glosarium}
\label{app:glosarium}

Istilah di bawah dipakai dengan arti teknis yang tetap sepanjang buku. Padanan
Inggris dicantumkan karena sebagian besar literatur memakai istilah Inggris.

\begin{description}[style=nextline,itemsep=0.35em]
  \item[Adjoin] matriks kofaktor yang disusun transpos; dipakai menuliskan invers sebagai satu rumus, $A^{-1}=\operatorname{adj}(A)/\det A$ (\emph{adjugate}).
  \item[Adjoint] dalam turunan balik, bilangan yang menyatakan seberapa besar kerugian akhir terpengaruh oleh satu besaran di tengah jaringan; satu adjoint per besaran (\emph{adjoint}).
  \item[Akar primitif] anggota yang bila dipangkatkan berulang melewati semua anggota grup siklik lain sebelum kembali ke satu; pada kriptografi, posisi awal yang mencapai seluruh posisi kunci (\emph{primitive root}).
  \item[Aksioma] pernyataan yang diterima tanpa bukti sebagai titik awal; seluruh isi bab bergantung padanya (\emph{axiom}).
  \item[Aksioma kelengkapan] janji bahwa setiap himpunan tak kosong yang terbatas ke atas di $\mathbb{R}$ punya supremum; inilah yang membuat garis bilangan tidak berlubang (\emph{completeness axiom}).
  \item[Amplitudo] amplitudo peluang: bilangan kompleks yang modulus kuadratnya memberi peluang; berbeda dari peluang karena dua amplitudo dapat saling meniadakan (\emph{probability amplitude}).
  \item[Anggaran privasi] jumlah $\varepsilon$ yang boleh dibelanjakan sebuah sistem; setiap kueri atau langkah pelatihan mengurangi sisanya (\emph{privacy budget}).
  \item[Antiturunan] fungsi yang turunannya adalah fungsi yang sedang diintegralkan; arah pulang dari sebuah turunan (\emph{antiderivative}).
  \item[Aturan Born] aturan yang mengubah amplitudo menjadi peluang dengan mengkuadratkan modulusnya (\emph{Born rule}).
  \item[Aturan rantai] laju perubahan fungsi tersusun sama dengan hasil kali laju pada setiap mata rantainya (\emph{chain rule}).
  \item[Aturan tiga] bila tidak ada kejadian dalam $n$ percobaan, laju kejadiannya masih boleh setinggi sekitar $3/n$ pada selang kepercayaan 95 persen (\emph{rule of three}).
  \item[Bantalan] isi tambahan pada paket sebelum dikirim agar setiap paket berukuran tetap, sehingga ukuran paket tidak lagi membocorkan isinya (\emph{padding}).
  \item[Barisan] deretan bilangan yang diberi nomor urut $a_1,a_2,\dots$; urutannya bagian dari definisinya, bukan sesuatu yang ditambahkan kemudian (\emph{sequence}).
  \item[Barisan Cauchy] barisan yang sukunya makin rapat satu sama lain, tanpa perlu menyebut nilai yang dituju (\emph{Cauchy sequence}).
  \item[Barisan monoton] barisan yang hanya naik atau hanya turun, tidak pernah berbalik arah (\emph{monotone sequence}).
  \item[Basis] himpunan vektor yang bebas linear dan merentang seluruh ruang (\emph{basis}).
  \item[Bentuk baku] penulisan seragam masalah optimasi: satu sasaran, kendala berbentuk tidak lebih dari nol, dan peubah tak negatif (\emph{standard form}).
  \item[Bias-variansi] pembagian galat model menjadi galat karena terlalu sederhana dan galat karena terlalu peka pada data (\emph{bias-variance tradeoff}).
  \item[Bijeksi] pemasangan sempurna dua himpunan: injektif dan surjektif sekaligus, tanpa sisa di kedua sisi (\emph{bijection}).
  \item[Bilangan kondisi] ukuran seberapa besar galat relatif keluaran dapat membesar dibanding galat relatif masukan (\emph{condition number}).
  \item[Bootstrap] penaksiran sebaran statistik dengan mengambil ulang sampel dari data sendiri, berkali-kali, dengan pengembalian (\emph{bootstrap}).
  \item[BPDA] pengganti gradien pada potongan yang tidak terdiferensiasi dengan gradien fungsi di titik sesudahnya, dipakai agar serangan tetap dapat dihitung (\emph{backward pass differentiable approximation}).
  \item[Celah integralitas] jarak antara nilai terbaik relaksasi LP dan nilai terbaik jawaban bulat; ia mengukur seberapa jauh model bersambung menipu (\emph{integrality gap}).
  \item[Defective] matriks yang nilai eigennya berulang tetapi ruang eigennya terlalu sempit untuk menjadi basis (\emph{defective}).
  \item[Dekomposisi nilai singular (SVD)] penulisan matriks sebagai $U\Sigma V^{\top}$ dengan $U,V$ ortogonal (\emph{singular value decomposition}).
  \item[Deret] penjumlahan suku-suku barisan tanpa henti, yang nilainya dibaca sebagai limit jumlah parsial (\emph{series}).
  \item[Deret harmonik] deret $1+1/2+1/3+\cdots$; sukunya menuju nol tetapi jumlahnya menyimpang (\emph{harmonic series}).
  \item[Deret pangkat] deret berbentuk $\sum a_n(x-c)^n$; ia berlaku hanya di dalam jari-jari konvergensinya (\emph{power series}).
  \item[Deret Taylor] penulisan fungsi sebagai jumlah suku turunan di satu titik, beserta sisa yang mengukur galat pemotongan. (\emph{Taylor series}).
  \item[Determinan] satu bilangan yang mengukur faktor perubahan volume dan sekaligus menguji invertibilitas (\emph{determinant}).
  \item[Determinan Jacobian] faktor skala volume lokal ketika koordinat diganti; nilai mutlaknya mengalikan luas atau volume pada integral lipat (\emph{Jacobian determinant}).
  \item[Diagonalisasi] penggantian basis sampai matriksnya menjadi diagonal, sehingga pekerjaannya tinggal perkalian angka (\emph{diagonalization}).
  \item[Dimensi] banyaknya vektor dalam sebuah basis; angkanya milik ruang, bukan milik satu basis yang kebetulan dipilih (\emph{dimension}).
  \item[Distribusi eksponensial] sebaran waktu tunggu pada proses tanpa penuaan, dengan fungsi kepadatan $\lambda e^{-\lambda t}$ (\emph{exponential distribution}).
  \item[Distribusi normal] sebaran lonceng simetris yang ditentukan sepenuhnya oleh rata-rata dan variansinya (\emph{normal distribution}).
  \item[Dualitas] perumusan masalah optimasi dari sisi lain, yang nilai optimalnya memberi batas bawah bagi masalah asal (\emph{duality}).
  \item[Ekor berat] sebaran yang peluang nilai raksasanya mengecil terlalu lambat, sehingga satu pengamatan dapat mengalahkan seluruh sisanya (\emph{heavy tail}).
  \item[Eliminasi] penggantian sebuah sistem dengan sistem yang lebih sederhana tanpa mengubah himpunan solusinya (\emph{elimination}).
  \item[Entropi silang] ukuran kejutan $-\sum_i y_i\log \hat{p}_i$; makin yakin model pada jawaban yang salah, makin besar dendanya (\emph{cross-entropy}).
  \item[EOT] serangan yang memaksimumkan harapan kerugian atas keacakan pertahanan, biasanya dengan merata-ratakan gradien pada beberapa salinan acak (\emph{expectation over transformation}).
  \item[Epsilon mesin] jarak antara $1$ dan bilangan mesin berikutnya; ukuran satu langkah pada skala satu (\emph{machine epsilon}).
  \item[Epsilon-delta] gaya mendefinisikan limit dengan menyebut toleransi keluaran lebih dulu, lalu batas masukan yang menjaminnya. (\emph{epsilon-delta definition}).
  \item[Faktor integrasi] pengali yang mengubah ruas kiri persamaan diferensial linear menjadi satu turunan hasil kali (\emph{integrating factor}).
  \item[Fungsi aktivasi] pemetaan tak linear pada keluaran sebuah unit; ia menentukan seberapa besar sinyal diteruskan dan bentuk turunannya (\emph{activation function}).
  \item[Fungsi distribusi kumulatif] fungsi $F(x)=P(X\le x)$; tangga menanjak yang tingginya menyatakan peluang yang sudah terkumpul sampai $x$ (\emph{cumulative distribution function}).
  \item[Fungsi kerugian] bilangan yang ingin diperkecil saat model dilatih; ia mengukur seberapa salah prediksi, bukan seberapa benar (\emph{loss function}).
  \item[Fungsi konveks] fungsi yang grafiknya selalu di bawah tali penghubung dua titiknya; setiap minimum lokalnya adalah minimum global (\emph{convex function}).
  \item[Fungsi Lyapunov] fungsi bernilai positif yang menurun sepanjang lintasan sistem; keberadaannya membuktikan kestabilan tanpa menyelesaikan persamaannya (\emph{Lyapunov function}).
  \item[Fungsi phi] fungsi $\varphi(n)$: banyaknya bilangan antara $1$ dan $n$ yang saling prima dengan $n$ (\emph{Euler totient function}).
  \item[Galat baku] simpangan baku penaksir, bukan simpangan baku data; ia mengukur goyangan satu tebakan (\emph{standard error}).
  \item[Galat pemotongan] selisih yang muncul karena proses tak hingga diganti potongan berhingga (\emph{truncation error}).
  \item[Garis regresi] garis (atau bidang) yang meminimumkan jumlah kuadrat sisa; ekuivalen dengan proyeksi ortogonal pada ruang kolom (\emph{regression line}).
  \item[Generalisasi] kemampuan mengerjakan soal yang belum pernah dilihat, bukan soal yang sudah dihafal (\emph{generalization}).
  \item[Gradien] vektor turunan parsial; arah naik tercepat sebuah fungsi (\emph{gradient}).
  \item[Gradien dikaburkan] gradien yang diberikan kepada penyerang tidak menunjuk arah yang benar, biasanya karena ada potongan tak terdiferensiasi atau acak (\emph{obfuscated gradient}).
  \item[Gradien terbalik] lihat turunan balik. (\emph{backward gradient}).
  \item[Harapan] nilai harapan $\mathbb{E}[X]$: rata-rata berbobot peluang, titik tumpu sebaran (\emph{expectation}).
  \item[Hasilkali dalam] pemetaan dua vektor ke satu skalar yang mengukur panjang dan sudut; di $\mathbb{R}^n$ nilainya $\langle \mathbf{x},\mathbf{y}\rangle=\mathbf{x}^{\top}\mathbf{y}$ (\emph{inner product}).
  \item[Hessian] matriks turunan parsial kedua; kecekungan fungsi ditentukan oleh tandanya (\emph{Hessian}).
  \item[Himpunan kuasa] himpunan semua himpunan bagian, termasuk himpunan kosong dan himpunan itu sendiri; kardinalitasnya $2^n$ (\emph{power set}).
  \item[Himpunan padat] himpunan yang setiap titiknya dapat didekati oleh anggota himpunan itu sendiri selain titik tersebut (\emph{dense set}).
  \item[Himpunan terhitung] himpunan yang anggotanya dapat didaftar sebagai barisan pertama, kedua, ketiga, dan seterusnya tanpa ada yang tertinggal (\emph{countable set}).
  \item[Hipotesis nol] dugaan awal yang diuji, biasanya "tidak ada efek"; penolakan hanya terjadi bila data sangat tidak cocok dengannya (\emph{null hypothesis}).
  \item[Hukum bilangan besar] janji bahwa rata-rata banyak pengamatan menempel ke rata-rata sebenarnya, dengan peluang meleset yang mengecil (\emph{law of large numbers}).
  \item[i.i.d.] sampel yang saling bebas dan berdistribusi sama; asumsi paling sering dipakai, dan paling sering dilanggar (\emph{independent and identically distributed}).
  \item[Implikasi] pernyataan berbentuk jika $P$ maka $Q$; ia hanya salah bila $P$ benar dan $Q$ salah (\emph{implication}).
  \item[Independen] dua kejadian yang tidak saling memberi kabar, $P(A\cap B)=P(A)P(B)$; mengetahui yang satu tidak mengubah peluang yang lain (\emph{independent}).
  \item[Induksi] pola bukti domino: basis benar, lalu setiap kasus menjatuhkan kasus berikutnya (\emph{induction}).
  \item[Induksi kuat] induksi yang boleh memakai seluruh kasus sebelumnya, bukan hanya kasus terdekat (\emph{strong induction}).
  \item[Integral lipat dua] penjumlahan nilai pada petak-petak kecil sebuah daerah, lalu limitnya ketika petak mengecil (\emph{double integral}).
  \item[Integral Riemann] limit jumlah luas persegi panjang ketika lebar partisi mengecil (\emph{Riemann integral}).
  \item[Interpolasi] mencari fungsi yang melewati semua titik data; berbeda dari regresi yang tidak menuntut itu. (\emph{interpolation}).
  \item[Interval kepercayaan] rentang nilai yang, bila prosedur pengambilan sampel diulang, memuat parameter sebenarnya dengan peluang tertentu (\emph{confidence interval}).
  \item[Interval Wilson] selang kepercayaan proporsi yang dicari dengan membalik uji, sehingga tetap masuk akal ketika proporsi mendekati nol atau satu (\emph{Wilson interval}).
  \item[Inti] lihat kernel: himpunan semua vektor yang dipetakan ke nol oleh transformasi linear (\emph{kernel, null space}).
  \item[Invers modular] bilangan $b$ sehingga $ab\equiv 1 \pmod n$; tombol batalkan pada aritmetika jam (\emph{modular inverse}).
  \item[Isomorfisme] pemetaan bijektif yang mempertahankan struktur; dua ruang yang isomorfik tidak dapat dibedakan secara aljabar (\emph{isomorphism}).
  \item[Jacobian] matriks turunan parsial pertama sebuah pemetaan vektor; determinannya mengukur faktor skala volume lokal (\emph{Jacobian}).
  \item[Jari-jari konvergensi] jarak terjauh dari pusat tempat sebuah deret pangkat masih konvergen (\emph{radius of convergence}).
  \item[Jumlah atas] taksiran integral yang sengaja dibesarkan, memakai nilai terbesar pada tiap potongan (\emph{upper sum}).
  \item[Jumlah langsung] penjumlahan dua subruang yang hanya bertemu di vektor nol, sehingga setiap vektor punya satu cara penulisan (\emph{direct sum}).
  \item[Kalibrasi] kesesuaian antara angka yang dijanjikan model dan frekuensi kenyataan pada kelompok yang dijanjikan sama (\emph{calibration}).
  \item[Kamus] kotak kecil di dalam bab yang menerangkan satu istilah baru beserta analoginya.
  \item[Kebocoran data] informasi dari data uji yang tanpa sengaja masuk ke proses pelatihan, sehingga hasil uji tampak lebih baik dari kenyataan (\emph{data leakage}).
  \item[Kekontinuan seragam] kekontinuan dengan satu $\delta$ yang berlaku untuk semua titik pada himpunan, bukan satu $\delta$ per titik (\emph{uniform continuity}).
  \item[Kekonvergenan seragam] barisan fungsi yang mendekati fungsi batasnya dengan laju yang sama di seluruh domain (\emph{uniform convergence}).
  \item[Kelas sisa] himpunan bilangan yang memberi sisa sama saat dibagi $n$, seperti pukul tiga pagi dan pukul tiga sore (\emph{residue class}).
  \item[Kernel] himpunan semua masukan yang dipetakan ke nol oleh transformasi linear; sama dengan ruang nol (\emph{kernel, null space}).
  \item[Kesalahan pembulatan] selisih antara bilangan real dan wakilannya di mesin, yang menumpuk pada penjumlahan panjang (\emph{rounding error}).
  \item[Kestabilan] sifat titik setimbang atau sistem dinamik yang tetap dekat dengan keadaannya setelah gangguan kecil (\emph{stability}).
  \item[Keterintegralan] sifat fungsi yang jumlah atas dan bawahnya dapat dirapatkan sampai selisihnya hilang, disebut keterintegralan Riemann (\emph{Riemann integrability}).
  \item[Ketunggalan] jaminan bahwa satu keadaan awal hanya melahirkan satu lintasan solusi (\emph{uniqueness}).
  \item[Koefisien determinasi] bagian variasi yang dijelaskan model, dinyatakan sebagai pecahan dari keseluruhan; ditulis $R^2$ (\emph{coefficient of determination}).
  \item[Kombinasi] pilihan beberapa anggota tanpa memperhatikan urutannya; banyaknya $\binom{n}{k}$ (\emph{combination}).
  \item[Kombinasi konveks] campuran berbobot vektor dengan bobot tak negatif yang berjumlah satu (\emph{convex combination}).
  \item[Kompak] himpunan tertutup dan terbatas; setiap barisan di dalamnya punya subbarisan yang konvergen di dalam himpunan itu (\emph{compact}).
  \item[Kompleksitas sampel] banyaknya contoh yang dibutuhkan agar model dapat digeneralisasi dengan keyakinan tertentu (\emph{sample complexity}).
  \item[Kondisi KKT] syarat perlu optimalitas untuk masalah dengan kendala pertidaksamaan (\emph{Karush-Kuhn-Tucker conditions}).
  \item[Kontinu] fungsi yang grafiknya dapat digambar tanpa mengangkat pena di titik itu, yaitu $\lim_{x\to c}f(x)=f(c)$ (\emph{continuous}).
  \item[Konvergensi dalam distribusi] fungsi distribusinya menuju fungsi distribusi batas di setiap titik kontinuitas; yang menyerupai adalah bentuk sebarannya, bukan nilainya (\emph{convergence in distribution}).
  \item[Konvergensi dalam peluang] $P(|X_n-X|>\varepsilon)\to 0$ untuk setiap $\varepsilon>0$ (\emph{convergence in probability}).
  \item[Konvergensi mutlak] deret yang tetap konvergen ketika semua sukunya dipaksa menjadi positif (\emph{absolute convergence}).
  \item[Konvergensi seragam] lihat kekonvergenan seragam: satu laju untuk seluruh domain, bukan satu jadwal per titik (\emph{uniform convergence}).
  \item[Koordinat] daftar koefisien yang menyatakan sebuah vektor dalam suatu basis; disebut juga vektor koordinat (\emph{coordinate vector}).
  \item[Korelasi] ukuran kekuatan hubungan linear antara dua peubah; tidak menyatakan sebab-akibat (\emph{correlation}).
  \item[Kovarians] lihat kovariansi: ukuran dua besaran bergerak bersama, dihitung dari hasil kali simpangan keduanya (\emph{covariance}).
  \item[Kovariansi] ukuran seberapa dua peubah bergerak bersama; satuannya hasil kali satuan keduanya (\emph{covariance}).
  \item[Kuadrat terkecil] pemilihan parameter yang meminimumkan jumlah kuadrat sisa ketika tidak ada solusi eksak (\emph{least squares}).
  \item[Kuantisasi] penggantian setiap nilai dengan kelipatan terdekat dari jarak tetap; hasilnya fungsi tangga yang turunannya nol hampir di mana-mana (\emph{quantization}).
  \item[Kuantor] penanda untuk semua ($\forall$) atau ada ($\exists$); urutan pemakaiannya mengubah arti pernyataan (\emph{quantifier}).
  \item[Kurva ROC] grafik tingkat benar terhadap tingkat salah palsu pada semua ambang; luas di bawahnya disebut AUROC (\emph{ROC curve}).
  \item[Kutub dan nol] nilai peubah kompleks yang membuat fungsi meledak atau lenyap; dipakai dalam analisis sistem dinamik dan kontrol. (\emph{poles and zeros}).
  \item[Lapangan] himpunan dengan dua operasi yang tertutup, punya identitas dan invers, dan memenuhi hukum distributif (\emph{field}).
  \item[Lapangan terurut] lapangan yang anggotanya dapat diurutkan dan aturan hitungnya menghormati urutan itu (\emph{ordered field}).
  \item[Lapisan tersembunyi] lapisan unit yang tidak bersentuhan langsung dengan masukan maupun keluaran jaringan (\emph{hidden layer}).
  \item[Lawan contoh] satu contoh yang mematahkan pernyataan "untuk semua" (\emph{counterexample}).
  \item[Limit fungsi] $f(x)$ dapat dipaksa sedekat apa pun ke $L$ asalkan $x$ cukup dekat ke $c$ tetapi tidak sama dengan $c$ (\emph{functional limit}).
  \item[Limit superior] nilai terbesar yang masih dikunjungi barisan sesering apa pun di ujungnya, ditulis $\limsup_{n\to\infty}a_n$ (\emph{upper limit}).
  \item[Lipschitz] syarat bahwa laju perubahan sebuah fungsi dibatasi satu konstanta; inilah yang menjamin solusi persamaan diferensial tunggal (\emph{Lipschitz condition}).
  \item[Luar sampel] data yang tidak dipakai saat melatih model, sehingga jujur untuk menilai (\emph{out-of-sample}).
  \item[Matriks elementer] matriks identitas yang sudah dikenai satu operasi baris, sehingga menjadi catatan operasi itu (\emph{elementary matrix}).
  \item[Matriks konfusi] tabel empat sel yang membandingkan prediksi dengan kenyataan (\emph{confusion matrix}).
  \item[Matriks ortogonal] matriks persegi yang kolom-kolomnya ortonormal; ia mempertahankan panjang dan sudut (\emph{orthogonal matrix}).
  \item[Metode Euler] penyelesaian numerik persamaan diferensial dengan melangkah mengikuti garis singgung; galatnya orde satu terhadap langkah. (\emph{Euler method}).
  \item[Metode tanda gradien] serangan satu langkah (FGSM) yang menggeser setiap koordinat sebesar anggaran, arahnya mengikuti tanda turunan kerugian (\emph{fast gradient sign method}).
  \item[Model ancaman] daftar lengkap kemampuan penyerang: besar gangguannya, apa yang ia ketahui, dan apa tujuannya (\emph{threat model}).
  \item[Multikolinearitas] keadaan ketika peubah penjelas saling hampir bergantung linear, sehingga koefisien regresi menjadi tidak stabil (\emph{multicollinearity}).
  \item[Multiplisitas] banyaknya kemunculan sebuah nilai eigen sebagai akar polinomial karakteristik, yang tidak selalu sama dengan dimensi ruang eigennya (\emph{multiplicity}).
  \item[NIDS] sistem yang memantau lalu lintas jaringan dan memutuskan apakah suatu aliran paket merupakan serangan (\emph{network intrusion detection system}).
  \item[Nilai eigen] skalar $\lambda$ sehingga $A\mathbf{v}=\lambda\mathbf{v}$ untuk suatu vektor tak nol (\emph{eigenvalue}).
  \item[Nilai mutlak] jarak sebuah bilangan ke nol pada garis bilangan, tanpa memedulikan tandanya (\emph{absolute value}).
  \item[Nilai singular] akar nilai eigen dari $A^{\top}A$; mengukur panjang sumbu utama pemetaan oleh matriks $A$ (\emph{singular value}).
  \item[Nilai-p] peluang memperoleh hasil setidaknya seekstrem yang teramati bila hipotesis nol benar; bukan peluang hipotesis nol benar (\emph{p-value}).
  \item[Norma] ukuran panjang vektor; norma $L_\infty$ mengambil komponen terbesar. (\emph{norm}).
  \item[Orientasi] kesepakatan tentang arah putar mana yang dihitung positif; membaliknya membalik tanda integral (\emph{orientation}).
  \item[Ortogonal] saling tegak lurus dalam arti hasilkali dalamnya nol (\emph{orthogonal}).
  \item[Ortonormal] himpunan ortogonal yang setiap anggotanya berpanjang satu (\emph{orthonormal}).
  \item[Overfitting] menghafal data latih termasuk deraunya, sehingga kemampuan di data baru menurun (\emph{overfitting}).
  \item[Partisi] pembagian himpunan menjadi bagian-bagian yang saling lepas dan menutup seluruh himpunan (\emph{partition}).
  \item[Peluang] bobot pada setiap kemungkinan, tak negatif dan berjumlah satu (\emph{probability}).
  \item[Peluang bersyarat] peluang kejadian bila kejadian lain diketahui terjadi; ditulis $P(A\mid B)$ (\emph{conditional probability}).
  \item[Pembatalan katastrofik] hilangnya digit berarti ketika dua bilangan hampir sama dikurangkan, sehingga yang tersisa hanya galat (\emph{catastrophic cancellation}).
  \item[Pembelajaran adversarial] studi tentang contoh yang sengaja dirancang untuk menipu model, dan cara mempertahankannya (\emph{adversarial learning}).
  \item[Pembulatan] penggantian bilangan real dengan bilangan mesin terdekat (\emph{rounding}).
  \item[Pemisahan data] pembagian data menjadi latih, validasi, dan uji yang tidak saling bocor (\emph{data split}).
  \item[Penaksir] aturan yang mengubah sampel menjadi tebakan parameter; tebakan yang dihasilkan disebut nilai dugaan (\emph{estimator}).
  \item[Pengali Lagrange] peubah bantu yang mengenakan harga pada kendala; nilainya menunjukkan seberapa besar tujuan dapat membaik bila kendala dilonggarkan (\emph{Lagrange multiplier}).
  \item[Penghubung logika] operasi yang menggabungkan pernyataan: dan, atau, tidak, jika-maka, jika-dan-hanya-jika (\emph{logical connective}).
  \item[Penurunan gradien stokastik] penurunan gradien yang memakai sebagian kecil data tiap langkah; lebih murah, lebih berisik (\emph{stochastic gradient descent}).
  \item[Perceptron] pengklasifikasi linear paling sederhana: tanda dari hasilkali dalam bobot dan masukan. (\emph{perceptron}).
  \item[Persamaan diferensial] kalimat yang mengikat laju perubahan sebuah besaran dengan besar besaran itu sendiri (\emph{differential equation}).
  \item[Persamaan diferensial biasa] persamaan yang mengaitkan fungsi dengan turunannya terhadap satu peubah bebas (\emph{ordinary differential equation}).
  \item[PGC (penurunan gradien)] aturan pembaruan parameter melawan arah gradien (\emph{gradient descent}).
  \item[Plateau tandus] daerah pada lanskap kerugian sirkuit variasional tempat gradien mengecil secara eksponensial terhadap jumlah qubit (\emph{barren plateau}).
  \item[Polinomial karakteristik] $\det(A-\lambda I)$; akar-akarnya adalah nilai eigen $A$ (\emph{characteristic polynomial}).
  \item[Privasi diferensial] janji bahwa keikutsertaan satu orang di dalam data hanya mengubah sebaran keluaran sedikit, diukur oleh $\varepsilon$ (\emph{differential privacy}).
  \item[Program linear] masalah memaksimumkan atau meminimumkan fungsi linear dengan kendala linear (\emph{linear program}).
  \item[Proposisi] kalimat yang bernilai benar atau salah, bukan pertanyaan atau perintah; dalam buku ini juga nama lingkungan teorema (\emph{proposition}).
  \item[Proyeksi] titik pada subruang yang paling dekat dengan sebuah vektor, yaitu bayangan tegak lurusnya (\emph{projection}).
  \item[Pseudoinvers] invers tergeneralisasi $A^{+}$ yang tetap bekerja ketika sebagian arah sudah gepeng: yang dapat dibalik dibalik, yang hilang dibiarkan (\emph{pseudoinverse}).
  \item[Rank] banyaknya arah bebas yang masih bertahan setelah dipetakan, yaitu $\dim\operatorname{im}T$ (\emph{rank}).
  \item[Rank-nullity] $\dim\ker T+\dim\operatorname{im}T=\dim V$ untuk transformasi linear pada ruang berdimensi hingga. (\emph{rank-nullity theorem}).
  \item[Rasio peluang] perbandingan peluang berhasil dan gagal; eksponensial koefisien regresi logistik menyatakan berapa kali rasio itu berlipat per satuan (\emph{odds ratio}).
  \item[Regresi] menarik fungsi yang paling dekat ke semua titik, bukan fungsi yang melewati semuanya (\emph{regression}).
  \item[Regresi logistik] model untuk keluaran biner yang memetakan kombinasi linear ke peluang lewat fungsi logistik (\emph{logistic regression}).
  \item[Regularisasi] penalti pada besarnya parameter agar model tidak terlalu peka pada data latih (\emph{regularization}).
  \item[Relaksasi LP] penggantian syarat bulat dengan syarat real agar masalah integer dapat diselesaikan sebagai LP (\emph{LP relaxation}).
  \item[Residu] selisih antara nilai teramati dan nilai yang diprediksi model (\emph{residual}).
  \item[Ringkasan] kotak butir di akhir bab yang boleh dibaca sendiri tanpa badan bab.
  \item[Ruang eigen] seluruh vektor eigen untuk satu nilai eigen beserta vektor nol; ia sebuah subruang (\emph{eigenspace}).
  \item[Ruang masalah] ruang objek nyata yang boleh diubah penyerang, misalnya paket dan waktu; berbeda dari ruang fitur tempat model bekerja (\emph{problem space}).
  \item[Ruang vektor] himpunan dengan penjumlahan dan perkalian skalar yang tertutup dan memenuhi delapan aksioma (\emph{vector space}).
  \item[Runge-Kutta] keluarga metode numerik yang mengevaluasi turunan beberapa kali per langkah untuk galat lebih kecil (\emph{Runge-Kutta}).
  \item[Selang prediksi] selang bagi satu pengamatan yang belum terjadi; selalu lebih lebar daripada selang bagi rata-rata (\emph{prediction interval}).
  \item[Sensitivitas] dalam privasi diferensial, perubahan keluaran terbesar yang dapat disebabkan satu orang (\emph{sensitivity}).
  \item[Sensitivitas global] pergeseran terbesar jawaban sebuah kueri yang dapat disebabkan satu orang pada kasus terburuk (\emph{global sensitivity}).
  \item[Serangan adaptif] serangan yang dirancang dengan mengetahui pertahanan yang sedang diuji, termasuk mengetahui bahwa gradiennya dimanipulasi (\emph{adaptive attack}).
  \item[Serangan adversarial] gangguan kecil pada masukan yang sengaja dibuat untuk mengubah keluaran model (\emph{adversarial example}).
  \item[Serupa] dua matriks yang menyatakan pemetaan linear yang sama dalam dua basis berbeda, $B=P^{-1}AP$ (\emph{similar}).
  \item[Sistem floating-point] notasi ilmiah biner dengan jumlah digit tetap; bilangan disimpan sebagai tanda, pangkat, dan mantissa, sehingga jarak antar bilangan yang dapat diwakili tidak seragam: rapat di dekat nol, renggang di ujung atas, dan dua bilangan berbeda dapat jatuh ke tempat yang sama (\emph{floating-point system}).
  \item[Softmax] pemetaan deretan skor menjadi peluang yang berjumlah satu, $e^{z_i}/\sum_j e^{z_j}$ (\emph{softmax}).
  \item[Solusi layak basis] titik yang diperoleh dengan menyetel sebagian peubah ke nol lalu menyelesaikan sisanya, dan masih memenuhi semua kendala (\emph{basic feasible solution}).
  \item[Solusi umum] keluarga semua solusi sebuah persamaan diferensial, bergantung pada konstanta sembarang (\emph{general solution}).
  \item[Sorry] penanda Lean~4 bahwa sebuah bukti ditangguhkan; mesin menerimanya tetapi mencatatnya sebagai aksioma (\emph{sorry}).
  \item[Stabil asimtotik] tetap dekat setelah diganggu, lalu benar-benar kembali ke titik setimbang (\emph{asymptotically stable}).
  \item[Subbarisan] barisan yang diperoleh dengan membuang sebagian suku tanpa mengubah urutan sisa sukunya (\emph{subsequence}).
  \item[Subruang] himpunan bagian ruang vektor yang tertutup terhadap penjumlahan dan perkalian skalar (\emph{subspace}).
  \item[Supremum] batas atas terkecil; bilangan terkecil yang tidak lebih kecil dari semua anggota himpunan (\emph{supremum, least upper bound}).
  \item[Susunan ulang] penukaran urutan suku deret; pada deret yang konvergen mutlak jumlahnya tidak berubah, pada deret yang tidak, jumlahnya dapat berubah (\emph{rearrangement}).
  \item[Syarat perlu] syarat yang harus dipenuhi agar pernyataan benar; tanpanya mustahil, tetapi memenuhinya belum cukup (\emph{necessary condition}).
  \item[Taktik] perintah langkah demi langkah dalam Lean~4 yang mengubah sasaran pembuktian (\emph{tactic}).
  \item[Tanda tangan digital] nilai yang hanya dapat dibuat pemegang kunci privat tetapi dapat diperiksa siapa pun dengan kunci publik (\emph{digital signature}).
  \item[Teorema Limit Pusat] rata-rata banyak variabel acak bebas menghampiri distribusi normal setelah dibakukan (\emph{central limit theorem}).
  \item[Terhitung] himpunan yang anggotanya dapat dinomori tanpa ada yang terlewat; $\mathbb{Q}$ terhitung, $\mathbb{R}$ tidak (\emph{countable}).
  \item[Term bukti] ungkapan yang sekaligus menjadi bukti dan menjadi nilai; memeriksa bukti sama dengan memeriksa tipe (\emph{proof term}).
  \item[Titik akumulasi] titik yang setiap persekitarannya memuat anggota himpunan lain (\emph{accumulation point}).
  \item[Titik setimbang] keadaan sistem dinamik yang tidak berubah bila tidak diganggu (\emph{equilibrium point}).
  \item[Transferabilitas] sifat contoh adversarial yang dibuat untuk satu model tetapi juga menipu model lain yang dilatih terpisah (\emph{transferability}).
  \item[Transformasi linear] pemetaan yang mempertahankan penjumlahan dan perkalian skalar (\emph{linear map}).
  \item[Turunan] limit hasil bagi selisih; kemiringan garis singgung sekaligus laju perubahan sesaat (\emph{derivative}).
  \item[Turunan balik] algoritme menghitung gradien seluruh parameter dengan menerapkan aturan rantai dari keluaran ke masukan (\emph{backpropagation}).
  \item[Typeclass] daftar kemampuan yang harus dimiliki sebuah tipe agar dapat dipakai di tempat tertentu (\emph{type class}).
  \item[Uji hipotesis] prosedur memutuskan, dengan risiko galat yang dinyatakan, apakah data cukup menolak dugaan awal (\emph{hypothesis test}).
  \item[Uji rasio] deret konvergen mutlak bila $\lim|a_{n+1}/a_n|<1$, dan divergen bila limit itu lebih besar dari satu (\emph{ratio test}).
  \item[ULP] satuan di tempat terakhir: satu langkah terkecil yang dapat dibedakan mesin pada skala sebuah bilangan (\emph{unit in the last place}).
  \item[Uniter] matriks yang mempertahankan hasilkali dalam, $U^{*}U=I$; versi kompleks dari matriks ortogonal (\emph{unitary}).
  \item[Validasi silang] penilaian model dengan membagi data bergiliran menjadi latih dan uji (\emph{cross-validation}).
  \item[Validasi silang bersarang] dua lapis lipatan: lapisan luar menilai, lapisan dalam memilih, sehingga pemilihan model tidak membocorkan data uji (\emph{nested cross-validation}).
  \item[Variabel acak] fungsi yang memetakan hasil percobaan ke bilangan (\emph{random variable}).
  \item[Variabel kendur] peubah tambahan yang mencatat sisa pemakaian, sehingga kendala tidak lebih dari berubah menjadi persamaan (\emph{slack variable}).
  \item[Varians inflasi] seberapa banyak variansi sebuah koefisien digelembungkan oleh penjelas lain yang membawa informasi serupa; disingkat VIF (\emph{variance inflation factor}).
  \item[Variansi] ukuran sebaran di sekitar rata-rata; akar kuadratnya adalah simpangan baku (\emph{variance}).
  \item[Verifikasi formal] pembuktian bahwa program atau rumus memenuhi spesifikasinya, diperiksa mesin (\emph{formal verification}).
  \item[Verteks] sudut daerah layak, yaitu titik yang tidak pernah menjadi titik tengah antara dua titik layak lain (\emph{vertex}).
  \item[Waktu polinomial] kelas algoritme yang banyaknya langkah tumbuh tidak lebih cepat dari polinom ukuran masukan (\emph{polynomial time}).
  \item[Wilson] interval kepercayaan proporsi yang tetap masuk akal ketika proporsi mendekati nol atau satu. (\emph{Wilson interval}).
  \item[Zaman kuantum] era ketika komputer kuantum cukup besar untuk menjalankan algoritme Shor pada kunci yang dipakai hari ini. (\emph{quantum computing era}).
\end{description}
